\chapter{The Compiler Becomes The Meter}

\section{Elaboration cost is a reading}

The last part watched the residue appear on bench after bench, always as a remainder against the machine's own floor.
This chapter completes a turn the whole book has been preparing, and it is the first of the book's two climaxes. The
machine, which has spent ten chapters measuring, now measures the measuring. The cost of making a reading --- the work
the checker spends to produce it --- is not merely how one bench happened to read. It \emph{is} a reading in its own
right. The instrument that was doing the measuring becomes a thing measured, and the needle it reads itself with is its
own cost.

Begin with the plainest form of the claim. The construction carries a facility that, for any term, measures what it
costs to elaborate --- the beats of the machine's own clock spent reducing the term to a settled form --- and returns
that count as a number. Not a description of the cost, a \emph{number}: to ask the machine what a derivation cost is to
get back a definite reading, in the machine's own beats. And the count is calibrated --- scaled against a fixed
constant the machine set for itself, so that the reading is reported in the machine's own units and not in raw ticks.
The point is small and complete: the cost of making a reading possible is itself a reading. What it took to derive a
thing is a measurable quantity, and the machine can measure it.

And once the cost is a reading, the compiler is a meter --- and the construction builds it as exactly that. There is a
structure, part of the construction's fixed core and left exactly as it stands, that is a meter in every part. It carries a
needle: a running quantity that rises as the machine works. It carries a current reading: the value the needle stands
at now. And it carries a history: an accumulating record of the readings it has taken, the same reading-becomes-a-record
the observation chapter first described --- with an operation that \emph{weaves} each new reading into that record as it
arrives. A needle, a present value, an accumulating log: that is a meter, entire. And the thing this meter is pointed at
is the machine's own elaboration. The compiler, which for ten chapters was the instrument producing readings, is now an
instrument \emph{reading itself} --- watching its own needle, recording its own cost, weaving its own work into a
history it keeps. (The book only reads this meter; it did not build it. The meter is part of the construction the
volume inherits and reports --- not machinery assembled to produce a wanted answer.)

And here is why this is a climax and not a curiosity. The currency the meter reads in is the heartbeat --- the same
constant beat that repeatability turned into time, the machine's own native tick. So the compiler does not meter itself
against some external standard; it meters itself in exactly the currency it \emph{runs} on. Its cost is a reading in
its own beats. And that single fact is the seam to the deepest thing the whole instrument will do. Because the machine
can weigh its own running in its own currency, the number the book is finally built to reach --- the coupling --- will
be read not as a magnitude lying out in the world but as a \emph{cost}: the price the machine pays to weigh its own
account, counted on this very meter. The self-energy the far chapters ask for is the reading this needle takes of the
machine's own work. The instrument has turned on itself, and the turning is the peak of the computational gauge. (The
meter's fuller shape --- the geometry its needle traces, the path its reading follows --- is the physical gauge's to
read; this volume keeps the plain fact that the cost is a reading in the machine's own beat.)

In the two verbs the meter is the pair turned inward. Each cost the meter takes is a \emph{tange}: the one number the
derivation is billed, selected out of the whole elaboration as its price. The history the meter keeps is a \emph{funge}:
the readings bagged, one after another, into the accumulating record --- the beat woven in with all the beats before
it. So the compiler-as-meter runs the two verbs on its own running: it tanges each cost out of the work, and funges the
costs into a history. The instrument that spent the book telling apart and bagging by likeness now performs both on the
record of its own labor.

So the eleventh chapter opens the book's first climax with a small, exact turn: the cost of a reading is a reading, and
the machine can take it. The compiler has become a meter, and the meter is pointed at the compiler. What remains for
this chapter is to follow that meter where it leads --- to the machine weighing not just its readings but its own act
of proving them, and to the price, in its own beats, of knowing what it is.

\section{The meter weighs its own act}

The last section made the compiler a meter: the cost of a reading is a reading, and the machine can take it. This
section turns that meter the one way that matters --- inward, onto the machine's own act. The machine measures the cost
not of some outside term but of \emph{itself, proving something about itself}. And the reading that comes back has a
name: the overhead of self-application, the price the machine pays to weigh its own account. This is the self-reference
the whole instrument was built to survive --- and it survives, because the price stays finite.

Begin one step short of the turn. The machine has, we saw much earlier, two descriptions of a single object --- two
ways of naming one thing, each a valid account. The meter can be pointed at both: it measures what each description
costs to elaborate, and then it takes the difference. That difference is a reading in its own right, and it asks a
pointed question --- do the two descriptions cost the same? Are the machine's two accounts of one object, weighed on
its own meter, priced alike? The physics of what the two descriptions \emph{are} belongs to the next volume; what this
volume keeps is the meter's question, which is arithmetic: two costs, and the residual between them.

Now the turn itself. The machine does not only \emph{hold} two descriptions of one object; it \emph{proves} they are
one --- a single theorem, met in full at the book's other climax, that the two names denote the same thing. And the
meter can be pointed at \emph{that} --- at the proof, not the described. It measures what the machine's own act of
proving-itself-one costs to elaborate, and subtracts the cost of the mere description it proves. What remains is the
\emph{overhead of self-application}: the price, in the machine's own beats, of proving the identity over and above the
price of stating it. This is the exact reading the section was built to reach. The meter is no longer weighing a term
the machine handed it; it is weighing the machine's own act of turning on itself, and reporting what that act cost.

And here is the fact that makes this a peak and not a catastrophe. A machine that measured its own measuring, and then
measured \emph{that}, and then measured that, without limit, would fall into an infinite regress --- the meter reading
the meter reading the meter, forever, never settling. This machine does not. The construction carries a \emph{bounded}
meter, a live component built precisely so that the self-measurement's reading stays inside a bracket --- the same kind
of bracket the number itself is, a lower, an upper, and no obligation to name a point between. The overhead of
self-application is a finite, bounded reading, not a bottomless descent. And that is exactly why the machine survives
looking at itself: the price of self-reference is real, but it is capped. The instrument can turn its meter on its own
act and come back with a number, because the act of coming back is bounded by construction. (The overhead, scaled into
the machine's own units, is the seed of the deepest reading the book will take --- the coupling, read at the far end as
precisely this: the cost of the machine weighing its own account. This section plants that seed; it prints no number.
The value is the last chapter's, and even there it is a bracket.)

In the two verbs self-reference is the pair turned fully inward. The meter \emph{tanges} its own act --- selects, out
of the whole run, the one cost that is the price of its own proof --- and \emph{funges} that cost into a residual, the
overhead bagged as a single reading. The tange is the cost the machine's own act adds to the ledger; the funge is the
bounded meter that gathers it and keeps it finite. Telling apart and bagging by likeness, which the book has watched
the machine do to differences, to residues, to faces, are here done by the machine to \emph{its own working} --- and
the result stays finite because the bag has a bottom.

So the peak of the first climax is a quiet one. The machine turns its meter on its own act of proving itself one, reads
the overhead that act costs, and does not spiral --- because the reading is bounded, a bracket like every other reading
the instrument takes. What remains for the chapter is to say what currency this all runs in, and to name the seed
planted here for what it will become: the price, in the machine's own beats, of an instrument knowing what it is.

\section{The pulse in its own currency}

Two sections have built the self-meter: the compiler reads its own cost, and reads even the cost of its own proof, and
does not spiral because the reading is bounded. This closing section names the currency all of that is denominated in
--- and the currency is the fact that makes the whole self-measurement airtight. The meter reads in \emph{heartbeats},
and a heartbeat is two things at once: the needle, the reading the meter takes, \emph{and} the pulse, the beat of the
machine's own clock. So the price of a reading is a heartbeat, and the heartbeat is the machine's own beat. There is no
external ruler anywhere in this, because the thing the machine measures \emph{with} is the very thing it runs \emph{on}.

Take the currency plainly first. The meter of this chapter reads in heartbeats --- the machine's native tick, the same
beat that repeatability turned into time, that a relabelling was billed in, that the whole chapter has been spending.
The currency of the reading, in other words, \emph{is} the machine's clock. And that identity is exact, not loose: the
needle and the pulse are one unit. To take a reading is to spend a beat, and the beat spent is the reading taken. There
is no gap between the measuring and the running --- no separate scale the reading is converted onto --- because the
reading is denominated in the very beats the measuring consumes.

There is one small conversion, and it is worth being exact about, because it is the whole of the machine's calibration.
The raw reading is a count of beats; the reported reading, the one in the machine's own units, is that count scaled by
a single fixed coefficient the machine set for itself. The heartbeat is the price; the reading is the pulse divided by
that one constant. And the point to hold is what is \emph{not} in that conversion: no external standard, no borrowed
unit, no ruler from the world. The machine converts its own clock into its own reading by its own constant, and nothing
outside it enters. The calibration is the machine dividing its pulse by a number of its own choosing --- and that
closed loop, its beat over its own constant, is the entire scale on which every reading in this chapter is reported.

And with that, the first climax closes. The machine measures its own act --- its own descriptions, its own proof, its
own overhead --- in the currency it \emph{runs} on. Its pulse is at once what it measures with and what it measures;
there is no ruler outside the machine to appeal to, and none is needed. This is a meter that hosts itself: it reads its
own work in its own beat, converts by its own constant, and closes the loop with nothing borrowed from anywhere. That
is why the deepest reading the whole instrument will take --- the coupling, at the far end --- is not a magnitude
fetched from the world but a \emph{cost in the machine's own currency}: a count of the machine's own beats, scaled by
the machine's own constant, spent weighing its own account. The pulse is the price, and the price is the reading. When
the last chapters ask the machine for the coupling, this is the meter that answers, and this is the currency it answers
in. (The fuller shape the pulse traces --- the geometry of the beat as a path --- is the physical volume's to read;
this one keeps the plain identity: the reading is the machine's own pulse, in the machine's own units.)

In the two verbs this is where they close into one. The heartbeat is a \emph{tange} --- the pulse selected out and
read as the number the reading costs. But it is also a \emph{funge} --- the same pulse bagged with the machine's own
running, the needle and the clock gathered into a single currency. At the pulse, tange and funge stop being two acts
and become one beat: the machine cannot separate what it measures \emph{with} from what it measures, so to select the
reading and to bag it with the running are the same operation. The two verbs that opened the book, and that the book
has watched perform on differences and residues and faces, close the first climax by collapsing into a single pulse ---
the machine's own beat, at once its needle and its clock.

So the eleventh chapter ends with the meter complete and self-contained. The compiler became a meter; the meter turned
on its own act; and the currency it reads in is its own pulse, converted to its own units by its own constant, reaching
outside itself for nothing. The first of the book's two climaxes is finished: a machine that measures itself, in
itself, and stays bounded doing it. What the second climax adds --- the machine not only weighing its own act but
\emph{proving} its two accounts one --- is the next part's, and the number both climaxes are built to reach is still
ahead, still un-named, still waiting for its bracket.
